\begin{definition}[The network adversary and the deployed composition]
  \label{def:aba-netadv}
  \lean{PLTS.ABA.Net.NetState}\lean{PLTS.ABA.Net.NetState.corrupt}\lean{PLTS.ABA.Net.NetStep}\lean{PLTS.ABA.Net.netAdv}\lean{PLTS.System.mapIdle}\lean{PLTS.ABA.Net.wccPull}\lean{PLTS.ABA.Net.wccLift}\lean{PLTS.ABA.Net.netPre}\lean{PLTS.ABA.Net.netGroup}\lean{PLTS.ABA.Net.netFlat}\leanok
  \uses{def:aba-procnet, def:aba-labels, def:aba-params, def:relabel, def:syncproduct, def:process-algebra}
  The state $\mathrm{NetState}(n)$ of the network adversary is the round-tagged pools $pool(r,j)$ of the stage messages process $j$ has multicast in stage $r$ (D5), the pools $dpool(j)$ of the DECIDED payloads $j$ has multicast (D12$'$), and the corrupted set $F \subseteq [n]$.
  Every pool is empty initially and so is $F$.
  Each process owns a DECIDED pool, so a corrupted process may hold both bits and have them delivered selectively (D12$'$).

  Corruption is one row.
  On $\mathrm{fail}(k)$ the adversary inserts $k$ into $F$ when $k \notin F$ and $|F| < f$, and changes nothing otherwise (D1).
  The corruption budget is that guard, and it is a guard of a component: nothing outside the adversary enforces it, and no restriction of the system and no hypothesis on its traces is needed to hold the adversary to $f$ corruptions.

  The automaton $\mathrm{netAdv}(P)$ has twenty rows, all Dirac.
  The send rows write a pool: $\mathrm{gsnd}$ pools $m$ under its sender, $\mathrm{callG}$ pools the $\langle\mathrm{INPUT},b\rangle$ that the call multicasts, $\mathrm{retWPub}$ pools the published payload (D10), and $\mathrm{dsnd}$ pools $b$ under the guard $b \notin dpool(j)$, the write-once condition of the DECIDED relay.
  The delivery rows carry the soundness of the authenticated network and consume nothing: $\mathrm{gdlv}$ requires $m \in pool(r,j)$ and $\mathrm{ddlv}$ requires $b \in dpool(j)$, so no message reaches a receiver that its named sender did not multicast.
  The row for $\mathrm{retABA}(id,b)$ requires $b \in dpool(id)$: a process returns only a payload it has itself published.
  On $\mathrm{callABA}$, $\mathrm{retG}$, $\mathrm{callW}$, $\mathrm{retW}$ and $\mathrm{gcallLoop}$ nothing is sent and the adversary stands.

  Seven rows read $F$, and they are the only rules of any component that do.
  Five are the authorisations of the Byzantine drives (D11): $\mathrm{byzCallG}$, $\mathrm{byzCallGLoop}$, $\mathrm{byzRetG}$, $\mathrm{byzCallW}$ and $\mathrm{byzRetW}$ each carry $k \in F$ and nothing else, the first of them additionally pooling the driven $\langle\mathrm{INPUT},b\rangle$.
  Two are the injections, and they are the only rows of any component labelled $\tau$ before hiding: $\mathrm{byzG}$ pools an arbitrary stage message under a corrupted sender at an arbitrary round (D5), and $\mathrm{byzD}$ pools either DECIDED bit under one (D12$'$), so equivocation at the DECIDED layer is $\mathrm{byzD}$ fired twice on one sender.

  The coin oracle $\mathrm{WCC.specFamily}(P)$, Transition System~3, is the third component and the only one whose transitions are not all Dirac.
  It speaks $\mathrm{Lab}(n)$, so it is read over the extended alphabet along a partial map $\mathrm{wccPull}$.
  A shared label is its own; the drive $\mathrm{byzCallW}(r,k)$ is the coin call $\mathrm{callW}(r,k)$, the drive $\mathrm{byzRetW}(r,k,b)$ the coin return $\mathrm{retW}(r,k,b)$, and the fused return $\mathrm{retWPub}(r,id,c,b)$ the coin return $\mathrm{retW}(r,id,c)$; every other rendezvous label has no image.
  The construction $\mathrm{mapIdle}$ reads a system over one alphabet as a system over another along such a map: a label with an image steps by that image's rows, and a label without one is a self-loop.
  The Byzantine handshake drives therefore reach the oracle over hidden labels while the network rendezvous leave it untouched.
  The oracle keeps its own copy of the corrupted set and its own budget-guarded corruption, and its resolution threshold reads that copy.

  The composition is
  \[
    \mathrm{netPre}(P) \;=\;
      \bigl(\mathrm{syncProduct}_{j \in [n]}\ \mathrm{ABAProcN}(P,j)\bigr)
        \parallel \bigl(\mathrm{netAdv}(P) \parallel \mathrm{wccLift}(P)\bigr),
  \]
  \[
    \mathrm{netGroup}(P) \;=\;
      \mathrm{relabel}\bigl(
        \mathrm{abstract}_{\mathrm{netEvtLabels}}(\mathrm{netPre}(P))\bigr),
  \]
  \[
    \mathrm{netFlat}(P) \;=\; \mathrm{abstract}_H\bigl(\mathrm{netGroup}(P)\bigr),
  \]
  where $\mathrm{wccLift}(P)$ is the oracle read along $\mathrm{wccPull}$ and $H$ is the sub-protocol API of Definition~\ref{def:aba-labels}.
  Hiding $\mathrm{netEvtLabels}$ discharges the extended alphabet, so $\mathrm{netGroup}(P)$ is read back along the left embedding and speaks exactly $\mathrm{Lab}(n)$.
  The state space of $\mathrm{netFlat}(P)$ is
  \[
    \bigl([n] \to \mathrm{ABANodeN}(n)\bigr) \times
      \bigl(\mathrm{NetState}(n) \times
        (\mathbb{N} \to \mathrm{WCC.SpecState})\bigr).
  \]
\end{definition}
